\pagestyle{empty}
\tight
\begin{tblr}{width=\textwidth,colspec={|Q[c,m,1.9cm]|X[l,m]|},hlines,vlines,hspan=minimal,row{1}={bg=headblue},rowsep=1pt,colsep=3pt}
\SetCell{c,m}\hfil\logo\hfil & \textbf{POLITEKNIK NEGERI MALANG}\par\textbf{JURUSAN TEKNOLOGI INFORMASI}\par\textbf{PROGRAM STUDI: D4 TEKNIK INFORMATIKA} \\
\SetCell[c=2]{c} \textbf{RENCANA PEMBELAJARAN SEMESTER (RPS)} & \\
\end{tblr}
\begin{longtblr}{width=\textwidth,colspec={|Q[l,m,1.9cm]|X[0.85,l,m]|X[1.45,l,m]|X[0.75,l,m]|X[1.15,l,m]|},hlines,vlines,hspan=minimal,rowsep=1pt,colsep=2pt,row{1,3}={bg=headgray,font=\bfseries}}
MATA KULIAH & KODE & BOBOT (SKS)/jam & SEMESTER & TGL. PENYUSUNAN \\
Agama & RTI252001 & 2 SKS/2 jam & 2 & 29 Januari 2026 \\
OTORISASI & Dosen Pengembang RPS & Koordinator RMK & \SetCell[c=2]{l} Ka PRODI &  \\
 & Astrifidha Rahma Amalia, S.Pd., M.Pd. & Astrifidha Rahma Amalia, S.Pd., M.Pd. & \SetCell[c=2]{l}{\vspace{8mm}\par Dr. Ely Setyo Astuti, S.T., M.T.} &  \\
\SetCell[r=9]{l} Capaian Pembelajaran (CP) & \SetCell[c=4]{l,bg=headgray,font=\bfseries} Capaian Pembelajaran Lulusan Program Studi (CPL-Prodi) &  &  &  \\
 & CPL01 & \SetCell[c=3]{l} Menginternalisasi ketakwaan kepada Tuhan Yang Maha Esa, menaati hukum, disiplin, dan menunjukkan profesionalisme melalui etika profesi, pembelajaran sepanjang hayat, serta respons terhadap isu sosial dan teknologi. &  &  \\
 & \SetCell[c=4]{l,bg=headgray,font=\bfseries} Capaian Pembelajaran mata kuliah (CPL-MK) &  &  &  \\
 & CPMK0102 & \SetCell[c=3]{l} Mampu menginternalisasi nilai ketakwaan kepada Tuhan YME, Pancasila, kewarganegaraan, dan menunjukkan pembelajaran sepanjang hayat melalui disiplin diri, tanggung jawab sosial, dan adaptasi karir profesional. &  &  \\
 & \SetCell[c=4]{l,bg=headgray,font=\bfseries} Kemampuan akhir tiap tahapan belajar (Sub-CPMK) &  &  &  \\
 & SCPMK0102-01001 & \SetCell[c=3]{l} Mahasiswa mampu menghargai perbedaan keyakinan dan agama dalam interaksi akademik dan sosial. &  &  \\
 & SCPMK0102-01002 & \SetCell[c=3]{l} Mahasiswa mampu bertindak jujur dalam pengerjaan tugas, ujian, dan interaksi akademik. &  &  \\
 & SCPMK0102-01003 & \SetCell[c=3]{l} Mahasiswa menunjukkan kedisiplinan dan tanggung jawab dalam menyelesaikan kewajiban akademik tepat waktu. &  &  \\
 & SCPMK0102-01004 & \SetCell[c=3]{l} Mahasiswa mampu merefleksikan nilai-nilai keagamaan dalam pengambilan keputusan etis di lingkungan kampus. &  &  \\
\SetCell[r=2]{l} Penilaian & \SetCell[c=4]{l} *Nilai CPMK didapat dari agregasi nilai SUB-CPMK yang dibebankan &  &  &  \\
 & \SetCell[c=4]{l}{\tiny\begin{tblr}{width=\linewidth,colspec={|Q[c,m,0.1\linewidth]|Q[c,m,0.16\linewidth]|X[l,m]|X[l,m]|X[l,m]|X[l,m]|X[l,m]|X[l,m]|Q[c,m,0.08\linewidth]|},hlines,vlines,hspan=minimal,rowsep=1pt,colsep=0.7pt,row{1,2}={bg=headgray,font=\bfseries}}
Kode CPMK & Kode SCPMK & \SetCell[c=6]{c} Bobot per Bentuk Penilaian & & & & & & Total Bobot \\
 & & Tugas & Kuis 1 & UTS & Kuis 2 & Case Method & UAS & \\
\SetCell[r=4]{c} CPMK0102 & SCPMK0102-01001 & & & & & 15\%: diskusi konsep manusia, agama, dan tauhid & & 15\% \\
 & SCPMK0102-01002 & 10\%: resume konsep manusia & 5\%: tes lisan bacaan dan hafalan Al-Qur'an & 10\%: materi minggu 1--7 & & 15\%: diskusi ekologi, akhlak, dan IPTEKS & & 40\% \\
 & SCPMK0102-01003 & & & & 5\%: video simulasi etos kerja & & & 5\% \\
 & SCPMK0102-01004 & & & & & 25\%: diskusi etos kerja, ekonomi syariah, dan masyarakat Islam & 15\%: materi minggu 9--16 & 40\% \\
\SetCell[c=8]{r}\textbf{Total Per Penilaian} & & & & & & & & \textbf{100\%} \\
\end{tblr}} &  &  &  \\
Diskripsi Singkat Mata Kuliah & \SetCell[c=4]{l} Mata kuliah wajib umum (MKU) yang mengkaji Pendahuluan, Konsep Manusia dan Agama, Konsep Tauhid, Manusia dalam Islam, Wawasan Ekologi dalam Islam, Aktualisasi Akhlak, IPTEKS dalam Islam, Pandangan Islam tentang Etos Kerja, Konsep Ekonomi Syariah, serta Pembentukan Masyarakat Islam melalui hakikat keluarga dan fungsi masjid. Mahasiswa diajak memahami ajaran Islam berdasarkan Al-Qur'an dan Hadis yang dihubungkan dengan kehidupan nyata era sekarang, sehingga PAI menjadi pegangan dalam pembentukan karakter yang baik secara individu maupun sosial. &  &  &  \\
Bahan Kajian / Materi Pembelajaran & \SetCell[c=4]{l}\begin{enumerate}\item Pendahuluan: visi pendidikan agama, deskripsi mata kuliah, pendekatan kuliah, tata tertib, dan sistem penilaian\item Manusia dan Agama: konsep agama dan Islam, agama kebutuhan manusia, dimensi ajaran Islam, metode memahami Islam, misi agama Islam, masa depan agama\item Konsep Tauhid: tauhid kebutuhan manusia, problem syirik, dampak tauhid\item Konsep Manusia: manusia dalam pandangan Islam, fungsi dan tugas hidup manusia, hakikat kehidupan dunia dan akhirat\item Wawasan Ekologi dalam Islam: hakikat alam semesta, arti dan sifat sunnatullah, manfaat alam semesta, Islam dan wawasan lingkungan\item Aktualisasi Akhlak: makna perjuangan Rasul, aktualisasi misi Rasul, aplikasi sifat-sifat Rasul dalam ilmu teknologi, ibadah dan pembentukan akhlak\item IPTEKS: hakikat ilmu dalam Islam, paradigma iptek, sumber ilmu pengetahuan, aplikasi pola dzikir dan pikir\item Etos Kerja: etos kerja dalam Islam, motivasi kerja dalam Islam, aktualisasi jihad dalam pembangunan\item Ekonomi: pengertian, prinsip, dan etika ekonomi syariah, pemberdayaan umat, manajemen zakat, wakaf, infak, dan sedekah\item Masyarakat Islam: fungsi keluarga dalam Islam, pembentukan keluarga sakinah, masjid sebagai pusat peradaban\end{enumerate} &  &  &  \\
\SetCell[r=2]{l} Pustaka & \SetCell[c=4]{l} \textbf{Utama:}\begin{enumerate}\item Al Qur'an Al Karim dan Qur'an Android.\item Fadloli, Sri Nurkudri, Abd. Chalim. 2018. Pendidikan Agama Islam Pada Perguruan Tinggi Umum. Malang: AM Publishing.\item Kemenristekdikti. 2016. Modul Pendidikan Agama Islam Untuk Perguruan Tinggi. Jakarta: Dirjen Belmawa Kemenristekdikti.\item Samiun Jazuli, Ahzami. 2014. Kehidupan Dalam Pandangan Al-Qur'an. Jakarta: Gema Insani.\item Abdullah, M. Amin. 2012. Islamic Studies di Perguruan Tinggi. Yogyakarta: Pustaka Pelajar.\end{enumerate} &  &  &  \\
 & \SetCell[c=4]{l} \textbf{Pendukung:}\begin{enumerate}\item Ibnu Katsir, Tafsir al-Qur'an Online.\item Djakfar, H. Muhammad. 2010. Teologi Ekonomi Membumikan Tita Langit di Ranah Bisnis. Malang: UIN Maliki Press.\item Umar, Nasaruddin. 2014. Islam Fungsional Revitalisasi \& Reaktualisasi Nilai Islam. Jakarta: Gramedia.\end{enumerate} &  &  &  \\
Media Pembelajaran & \SetCell[c=2]{l} \textbf{Software:} PPT dan modul. &  & \SetCell[c=2]{l} \textbf{Hardware:} Komputer, laptop, dan proyektor. &  \\
Nama Dosen Pengampu & \SetCell[c=4]{l}\begin{enumerate}\item Astrifidha Rahma Amalia, S.Pd., M.Pd.\end{enumerate} &  &  &  \\
Matakuliah Syarat & \SetCell[c=4]{l} - &  &  &  \\
\end{longtblr}
\begin{longtblr}{width=\textwidth,colspec={|Q[c,m,0.06\textwidth]|X[1.35,l,m]|X[1.08,l,m]|X[1.16,l,m]|X[0.7,l,m]|X[1.24,l,m]|X[1.06,l,m]|X[0.98,l,m]|Q[c,m,0.05\textwidth]|},hlines,vlines,hspan=minimal,rowhead=2,row{1,2}={bg=headgreen,font=\bfseries},rowsep=1pt,colsep=1.5pt}
Mgg.\par Ke & Kemampuan Akhir Yang Direncanakan (Sub-CP-MK) & Bahan kajian (Materi Pembelajaran) & Bentuk dan Metode Pembelajaran & Estimasi Waktu & Pengalaman Belajar Mahasiswa & Kriteria \& Bentuk Penilaian & Indikator Penilaian & Bobot Penilaian (\%) \\
(1) & (2) & (3) & (4) & (5) & (6) & (7) & (8) & (9) \\
1 & SCPMK0102-01001 & Pendahuluan: visi pendidikan agama, deskripsi mata kuliah, pendekatan kuliah, tata tertib, dan sistem penilaian. & Bentuk: kuliah tatap muka/daring/luring. Metode: ceramah, diskusi, studi kasus, dan pembelajaran berbasis masalah. & TM: 1 x (2 x 50'). & Mahasiswa memperoleh RPS, menyepakati kontrak perkuliahan dan tugas, serta melakukan tanya jawab dengan dosen. & Tugas kelompok: lembar observasi dosen. Bentuk penilaian: Case Method. & Partisipasi aktif dalam diskusi; keterbukaan terhadap pendapat lain; kesopanan dan etika selama diskusi. & 3.75 \\
2 & SCPMK0102-01001 & Manusia dan Agama: konsep agama dan Islam, agama kebutuhan manusia, dan dimensi ajaran Islam. & Bentuk: kuliah tatap muka/daring/luring. Metode: ceramah, diskusi, studi kasus, dan pembelajaran berbasis masalah. & TM: 1 x (2 x 50'). & Mahasiswa berdiskusi kelompok tentang konsep manusia dan agama, mengumpulkan makalah, serta presentasi lewat PPT. & Tugas kelompok: lembar observasi dosen. Bentuk penilaian: Case Method. & Partisipasi aktif; kemampuan menghubungkan konsep manusia dan agama dengan argumen relevan; komunikasi yang jelas dan profesional. & 3.75 \\
3 & SCPMK0102-01001 & Manusia dan Agama: metode memahami Islam, misi agama Islam, dan masa depan agama. & Bentuk: kuliah tatap muka/daring/luring. Metode: ceramah, diskusi, studi kasus, dan pembelajaran berbasis masalah. & TM: 1 x (2 x 50'). & Mahasiswa berdiskusi kelompok, melakukan tanya jawab terhadap peserta presentasi, dan menghubungkan materi dengan kehidupan sehari-hari. & Tugas kelompok: lembar observasi dosen. Bentuk penilaian: Case Method. & Partisipasi aktif; ketepatan argumen dan contoh; keterbukaan menerima pandangan lain. & 3.75 \\
4 & SCPMK0102-01001 & Konsep Tauhid: tauhid kebutuhan manusia, problem syirik, dan dampak tauhid. & Bentuk: kuliah tatap muka/daring/luring. Metode: ceramah, diskusi, studi kasus, dan pembelajaran berbasis masalah. & TM: 1 x (2 x 50'). & Mahasiswa berdiskusi kelompok tentang konsep tauhid, presentasi PPT, dan menghubungkan materi dengan konteks sosial dan moral. & Tugas kelompok: lembar observasi dosen. Bentuk penilaian: Case Method. & Kemampuan menghubungkan konsep manusia, agama, dan tauhid; penyampaian ide yang efektif; etika diskusi. & 3.75 \\
5 & SCPMK0102-01002 & Konsep Tauhid: praktik bacaan Al-Qur'an sesuai tajwid. & Bentuk: kuliah tatap muka/daring/luring. Metode: diskusi dan dialog, studi kasus, dan praktikum. & TM: 1 x (2 x 50'). & Mahasiswa dilatih menjaga fokus dan ketelitian saat membaca ayat Al-Qur'an serta mengintegrasikan nilai bacaan dalam kehidupan sehari-hari. & Kriteria: ketepatan bacaan dan tajwid. Bentuk penilaian: Kuis 1 (tes lisan). & Kesesuaian bacaan dengan aturan tajwid; kelancaran membaca tanpa terbata-bata. & 2.5 \\
6 & SCPMK0102-01002 & Konsep Tauhid: hafalan ayat Al-Qur'an dan penghayatan. & Bentuk: kuliah tatap muka/daring/luring. Metode: diskusi dan dialog, studi kasus, dan praktikum. & TM: 1 x (2 x 50'). & Mahasiswa diberi tantangan menghafal ayat Al-Qur'an dan materi ajaran agama serta belajar lebih teliti memahami teks. & Kriteria: ketepatan hafalan. Bentuk penilaian: Kuis 1 (tes lisan). & Ketepatan dan konsistensi hafalan ayat; penghayatan dalam membaca dan pengucapan. & 2.5 \\
7 & SCPMK0102-01002 & Konsep Manusia: manusia dalam pandangan Islam, fungsi dan tugas hidup manusia, hakikat kehidupan dunia dan akhirat. & Bentuk: kuliah tatap muka/daring/luring. Metode: diskusi dan dialog, studi kasus, dan resume. & TM: 1 x (2 x 50'). & Mahasiswa menyaring informasi relevan dari berbagai sumber, menyusun resume terstruktur dan ringkas, serta berpikir kritis terhadap sumber. & Kriteria: resume. Bentuk penilaian: Tugas (esai). & Cakupan konsep utama secara akurat; kejelasan struktur dan bahasa resume; keterkaitan materi dengan kehidupan sehari-hari dan orisinalitas. & 10 \\
8 & SCPMK0102-01001, SCPMK0102-01002 & UTS: materi minggu 1--7. & Ujian tulis, closed book. & Sesuai jadwal asesmen. & Mahasiswa mengerjakan soal ujian materi minggu 1--7 secara mandiri dan jujur. & Bentuk penilaian: UTS. & Ketepatan penjelasan konsep manusia, agama, dan tauhid; ketepatan penggunaan dalil; koherensi argumen. & 10 \\
9 & SCPMK0102-01002 & Wawasan Ekologi dalam Islam: hakikat alam semesta, sunnatullah, manfaat alam semesta, dan wawasan lingkungan. & Bentuk: kuliah tatap muka/daring/luring. Metode: ceramah, diskusi, studi kasus, dan pembelajaran berbasis masalah. & TM: 1 x (2 x 50'). & Mahasiswa memperoleh pengetahuan tentang wawasan ekologi, berdiskusi kelompok, dan melakukan tanya jawab dengan dosen. & Tugas kelompok: lembar observasi dosen. Bentuk penilaian: Case Method. & Partisipasi aktif; ketepatan menghubungkan konsep dengan argumen atau contoh relevan; etika diskusi. & 5 \\
10 & SCPMK0102-01002 & Aktualisasi Akhlak: makna perjuangan Rasul, aktualisasi misi Rasul, aplikasi sifat Rasul dalam ilmu teknologi, ibadah dan pembentukan akhlak. & Bentuk: kuliah tatap muka/daring/luring. Metode: ceramah, diskusi, studi kasus, dan pembelajaran berbasis masalah. & TM: 1 x (2 x 50'). & Mahasiswa berdiskusi kelompok tentang aktualisasi akhlak, mengumpulkan makalah, dan presentasi lewat PPT. & Tugas kelompok: lembar observasi dosen. Bentuk penilaian: Case Method. & Partisipasi aktif; kejelasan penyampaian ide; keterbukaan menerima pandangan lain. & 5 \\
11 & SCPMK0102-01002 & IPTEKS: hakikat ilmu dalam Islam, paradigma iptek, sumber ilmu pengetahuan, aplikasi pola dzikir dan pikir. & Bentuk: kuliah tatap muka/daring/luring. Metode: ceramah, diskusi, studi kasus, dan pembelajaran berbasis masalah. & TM: 1 x (2 x 50'). & Mahasiswa berdiskusi kelompok tentang IPTEKS dalam Islam, tanya jawab terhadap peserta presentasi, dan menghubungkan materi dengan kehidupan sehari-hari. & Tugas kelompok: lembar observasi dosen. Bentuk penilaian: Case Method. & Ketepatan argumen dan contoh; komunikasi profesional; kesopanan dan sikap hormat selama diskusi. & 5 \\
12 & SCPMK0102-01003 & Etos Kerja: etos kerja dalam Islam dan motivasi kerja dalam Islam. & Bentuk: kuliah tatap muka/daring/luring. Metode: ceramah, diskusi, studi kasus, dan pembelajaran berbasis masalah. & TM: 1 x (2 x 50'). & Mahasiswa memahami konsep etos kerja, membuat video simulasi beserta dalil Al-Qur'an dan Hadis, serta mengembangkan keterampilan perangkat lunak pembuatan video. & Tugas kelompok: video simulasi beserta dalil. Bentuk penilaian: Kuis 2 (rubrik penilaian video simulasi). & Kesesuaian isi video dengan pemahaman topik; kreativitas dan penggunaan teknologi; kesesuaian durasi dengan materi. & 5 \\
13 & SCPMK0102-01004 & Etos Kerja: aktualisasi jihad dalam pembangunan. & Bentuk: kuliah tatap muka/daring/luring. Metode: ceramah, diskusi, studi kasus, dan pembelajaran berbasis masalah. & TM: 1 x (2 x 50'). & Mahasiswa memperoleh pengetahuan tentang etos kerja, berdiskusi kelompok, dan melakukan tanya jawab dengan dosen. & Tugas kelompok: lembar observasi dosen. Bentuk penilaian: Case Method. & Partisipasi aktif; ketepatan argumen dan contoh relevan; etika diskusi. & 5 \\
14 & SCPMK0102-01004 & Ekonomi Syariah: pengertian, prinsip, dan etika ekonomi syariah. & Bentuk: kuliah tatap muka/daring/luring. Metode: ceramah, diskusi, studi kasus, dan pembelajaran berbasis masalah. & TM: 1 x (2 x 50'). & Mahasiswa berdiskusi kelompok tentang ekonomi syariah, mengumpulkan makalah, dan presentasi lewat PPT. & Tugas kelompok: lembar observasi dosen. Bentuk penilaian: Case Method. & Partisipasi aktif; kejelasan penyampaian ide; keterbukaan menerima pandangan lain. & 5 \\
15 & SCPMK0102-01004 & Ekonomi Syariah: pemberdayaan umat melalui ekonomi syariah, manajemen zakat, wakaf, infak, dan sedekah. & Bentuk: kuliah tatap muka/daring/luring. Metode: ceramah, diskusi, studi kasus, dan pembelajaran berbasis masalah. & TM: 1 x (2 x 50'). & Mahasiswa berdiskusi kelompok, melakukan tanya jawab terhadap peserta presentasi, dan menghubungkan materi dengan kehidupan sehari-hari. & Tugas kelompok: lembar observasi dosen. Bentuk penilaian: Case Method. & Ketepatan argumen dan contoh; komunikasi profesional; kesopanan dan etika diskusi. & 7.5 \\
16 & SCPMK0102-01004 & Masyarakat Islam: fungsi keluarga dalam Islam, pembentukan keluarga sakinah, masjid sebagai pusat peradaban. & Bentuk: kuliah tatap muka/daring/luring. Metode: ceramah, diskusi, studi kasus, dan pembelajaran berbasis masalah. & TM: 1 x (2 x 50'). & Mahasiswa berdiskusi kelompok tentang masyarakat Islam, presentasi PPT, dan merefleksikan materi dalam konteks sosial dan moral. & Tugas kelompok: lembar observasi dosen. Bentuk penilaian: Case Method. & Partisipasi aktif; refleksi nilai keagamaan pada kasus nyata; etika diskusi. & 7.5 \\
17 & SCPMK0102-01004 & UAS: materi minggu 9--16. & Ujian tulis, closed book. & Sesuai jadwal asesmen. & Mahasiswa mengerjakan soal ujian materi minggu 9--16 secara mandiri dan jujur. & Bentuk penilaian: UAS. & Pemahaman etos kerja, ekonomi syariah, dan masyarakat Islam; refleksi nilai keagamaan dalam keputusan etis; penggunaan dalil. & 15 \\
\end{longtblr}
